%!TEX program = pdflatex
\documentclass[12pt,a4paper]{article}
% Общие параметры оформления лабораторной работы.
\usepackage[T2A]{fontenc}
\usepackage[utf8]{inputenc}
\usepackage[russian]{babel}
\usepackage{amssymb,amsmath,mathtools}
\usepackage{graphicx}
\usepackage[table]{xcolor}
\usepackage{booktabs,tabularx,array}
\usepackage{enumitem}
\usepackage{needspace}
\usepackage{fancyhdr}
\usepackage{titlesec}
\usepackage{tcolorbox}
\tcbuselibrary{breakable,skins}
\usepackage[hidelinks]{hyperref}

% Размер наборной полосы и интервалы, согласованные с ЛР1 и ЛР2.
\setlength{\textwidth}{16cm}
\setlength{\textheight}{24cm}
\setlength{\oddsidemargin}{0mm}
\setlength{\evensidemargin}{0mm}
\setlength{\topmargin}{-22mm}
\setlength{\headheight}{15pt}
\setlength{\headsep}{7mm}
\setlength{\footskip}{11mm}
\setlength{\parindent}{0pt}
\setlength{\parskip}{2pt}
\tolerance=10000
\emergencystretch=1.2em
\clubpenalty=10000
\widowpenalty=10000
\displaywidowpenalty=10000
\setlength{\abovedisplayskip}{8pt plus 2pt minus 2pt}
\setlength{\belowdisplayskip}{8pt plus 2pt minus 2pt}
\setlength{\abovedisplayshortskip}{5pt}
\setlength{\belowdisplayshortskip}{6pt}
\setlist[enumerate]{leftmargin=2.7em,topsep=5pt,itemsep=4pt,parsep=0pt}
\setlist[itemize]{leftmargin=2.7em,topsep=4pt,itemsep=3pt,parsep=0pt}
\renewcommand{\arraystretch}{1.12}
\setlength{\tabcolsep}{5pt}
\titleformat{\section}{\normalfont\bfseries\large}{\thesection.}{.5em}{}
\titlespacing*{\section}{0pt}{14pt plus 2pt minus 2pt}{8pt}
\titleformat{\subsection}{\normalfont\bfseries\normalsize}{\thesubsection.}{.5em}{}
\titlespacing*{\subsection}{0pt}{12pt plus 2pt minus 2pt}{6pt}
\titleformat{name=\subsection,numberless}{\normalfont\bfseries\normalsize}{}{0pt}{}

% Колонтитулы и первая страница.
\newcommand{\labtitle}{Две стадии: от скрытого времени к прогнозу}
\newcommand{\labtag}{Лабораторная работа~№~3-4}
\newcommand{\itmoidentity}{\includegraphics[width=22mm]{../branding/itmo_logo.pdf}}
\pagestyle{fancy}
\fancyhf{}
\fancyhead[L]{\includegraphics[width=13mm]{../branding/itmo_logo.pdf}}
\fancyhead[R]{\small Теория вероятностей}
\fancyfoot[C]{\small\thepage}
\renewcommand{\headrulewidth}{0pt}
\renewcommand{\footrulewidth}{0pt}
\fancypagestyle{firstpage}{\fancyhf{}\renewcommand{\headrulewidth}{0pt}\fancyfoot[C]{\small\thepage}}
\newcommand{\dochead}[1]{%
  \noindent\begin{minipage}[c]{.77\textwidth}
   {\small Теория вероятностей\par}\vspace{4pt}
   {\bfseries\large \labtag\par}\vspace{2pt}
   {\bfseries\large \labtitle\par}
  \end{minipage}\hfill
  \begin{minipage}[c]{.20\textwidth}\raggedleft\itmoidentity\end{minipage}\par
  \vspace{10pt}
}

% Краткие примечания студенческой версии.
\newenvironment{studentnote}{\par\medskip\begingroup\small\noindent}{\par\endgroup\medskip}
\newcommand{\stageoutcome}[1]{\par\medskip{\small\textbf{Результат этапа.} #1}\par\medskip}

% Смысловые пометки для преподавателя.
\definecolor{explainbg}{RGB}{239,246,250}
\definecolor{solvebg}{RGB}{238,247,240}
\definecolor{commentbg}{RGB}{253,248,235}
\newtcolorbox{teacherbox}[1][]{enhanced,breakable,sharp corners,boxrule=0pt,frame hidden,colback=explainbg,left=8pt,right=8pt,top=6pt,bottom=6pt,before skip=6pt,after skip=6pt,#1}
\newtcolorbox{solutionbox}[1][]{enhanced,breakable,sharp corners,boxrule=0pt,frame hidden,colback=solvebg,left=8pt,right=8pt,top=6pt,bottom=6pt,before skip=6pt,after skip=6pt,#1}
\newtcolorbox{warningbox}[1][]{enhanced,breakable,sharp corners,boxrule=0pt,frame hidden,colback=commentbg,left=8pt,right=8pt,top=6pt,bottom=6pt,before skip=6pt,after skip=6pt,#1}
\hypersetup{pdfauthor={ИТМО},pdftitle={Лабораторная работа № 3-4. Две стадии}}


\begin{document}
\linespread{1.055}\selectfont
\setlength{\parskip}{3pt}
\setlist[enumerate]{itemsep=4pt,topsep=5pt}
\thispagestyle{firstpage}
\dochead{Исследование двух последовательных случайных стадий}

\section*{Цель работы}
Исследовать время прохождения двух последовательных случайных стадий: от первых измерений только полного времени до математической модели, прогноза риска задержки и проверки центральной предельной теоремы.

\section*{Что представляет собой установка}
Вы работаете с \textbf{учебным компьютерным стендом}, который воспроизводит последовательный процесс с двумя внутренними стадиями. В зависимости от назначенного направления за одной и той же схемой стоят разные физические ситуации:
\begin{itemize}[topsep=4pt,itemsep=3pt]
\item \textbf{Биоинженерия.} Одна молекула переходит из состояния A в промежуточное B, затем в C. Исследуется ожидание двух последовательных переходов.
\item \textbf{ФБИТ.} Тревожное сообщение проходит через два радиоприёмника с периодическими окнами связи. Оно ждёт доступного окна у первого, затем у второго приёмника. Это контролируемый опыт, где начальные фазы задаются независимо и заново перед каждым запуском.
\item \textbf{Физика и наноструктуры.} Один носитель заряда последовательно проходит через два туннельных перехода; между ними существует промежуточное зарядовое состояние B.
\end{itemize}
В каждом направлении исследуется \emph{одна} назначенная установка и один вариант её параметров. Времена генерируются математической моделью, а не снимаются с реального физического прибора.

\noindent\textbf{Что измеряется.} Событие A отмечает начало процесса, C -- завершение. Сначала известна только продолжительность $T$ между ними. Хотя внутри есть две стадии, журнал не сообщает, сколько заняла каждая. Дополнительный датчик B позволит в \emph{новых} запусках отдельно получить $X$ (от A до B) и $Y$ (от B до C), так что $T=X+Y$.

\noindent\textbf{Зачем это нужно.} Если полное время иногда оказывается слишком большим, недостаточно просто назвать более медленную стадию. Нужно построить закон суммы и по нему рассчитать, насколько изменение \emph{каждого} модуля снизит вероятность превышения допустимого порога.

\begin{studentnote}
\textbf{Главный исследовательский вопрос.} Можно ли по одному только $T$ понять вклад двух стадий? Какое новое измерение снимет неопределённость и позволит проверить количественный прогноз риска?
\end{studentnote}

\section*{Порядок исследования}
Работа устроена как последовательность из пяти этапов. На каждом этапе сначала зафиксируйте собственный ответ или прогноз, затем сравните его с результатами стенда:
\begin{enumerate}[topsep=4pt,itemsep=2pt]
\item Записать полные времена $T$ и объяснить, почему понадобился датчик B.
\item После B найти законы стадий, совместное распределение и обосновать предположения.
\item Получить условное распределение при известном $T$ и закон суммы через свёртку.
\item Рассчитать три риска $q_0,q_1,q_2$ и проверить эффект обоих вмешательств.
\item Для $n=100$ проверить прогноз ЦПТ на 300 независимых партиях.
\end{enumerate}
\noindent\textbf{Начало работы.} Откройте \texttt{start.html} по ссылке преподавателя, например \texttt{start.html?track=bio\&variant=03}. Проверьте направление и номер варианта в верхней части страницы. При смене любого из них измерения начинаются заново. Свои расчёты и выводы переносите в краткий отчёт; скриншоты всех экранов не нужны.

\section{Наблюдение и новый датчик}
\begin{enumerate}
\item Выполните несколько одиночных запусков, затем серию; получите не менее шести полных времён \(T\) без промежуточной метки. На ленте двигайте временную границу и объясните смысл доли \(\#\{T_i<t\}/n\).
\item До подключения дополнительного датчика ответьте: позволяют ли наблюдения только полного времени установить, \emph{какой именно модуль} задерживает дольше? Укажите, какой информации хватает и чего недостаёт. Аргументируйте ответ с помощью двух возможных объяснений одной и той же наблюдаемой записи.
\item До включения B зафиксируйте прогноз: что дополнительная метка может показать на новых временных трассах? Подключите датчик B и выполните хотя бы 12 новых проходов. Теперь по событиям A, B, C определяются \(X,Y,T\). Объясните, почему старые записи нельзя дополнить меткой B задним числом.
\end{enumerate}
\stageoutcome{Короткий аргумент о необходимости нового измерения; одна запись с $T=X+Y$. Снимки всех экранов не требуются.}

\Needspace{4\baselineskip}\section{Модель стадий}
После не менее 12 измерений с датчиком B изучите механизм своего направления и откройте характеристики стенда. Значения параметров считайте заданными свойствами учебной установки: оценивать их статистическими методами не нужно.
\begin{enumerate}[resume]
\item Объясните физическую причину случайной задержки и запишите законы \(X\) и \(Y\), единицы измерения и смысл параметров. Сформулируйте \emph{предположение о независимости} стадий, укажите одну ситуацию, где оно может нарушиться.
\item Найдите совместную плотность \(f_{X,Y}(x,y)\) и точно укажите область, где она положительна. Для установок E времена экспоненциальны; для U времена равномерны на промежутках, определяемых периодами.
\end{enumerate}
\begin{studentnote}\textbf{Важно для ФБИТ.} Независимые равномерные фазы двух приёмников задаются заново при каждом опыте. Для постоянно работающей сети с фиксированными фазами независимость не следует из названия устройства.\end{studentnote}

\Needspace{4\baselineskip}\section{Что скрывает известное полное время?}
\begin{enumerate}[resume]
\item Перейдите от \((X,Y)\) к \((X,T)\), где \(T=X+Y\). Запишите обратное преобразование и его якобиан, затем найдите \(f_{X,T}(x,t)\).
\item Получите условную плотность \(f_{X\mid T}(x\mid t)\) с \emph{правильной областью определения}. Для \(t=\mathbb E[T]\) рассчитайте одну вероятность:
\[
\mathbb P\!\left(X>\frac{T}{2}\,\middle|\,T=t\right).
\]
Объясните ответ человеческим языком и проверьте его в интерфейсе. При ошибке используйте подсказки и повторите расчёт; ответ можно раскрыть отдельно после нескольких попыток. \emph{Эмпирическая проверка условного закона на маленькой подвыборке не требуется.}
\end{enumerate}

\Needspace{4\baselineskip}\section{Свёртка и количественный прогноз риска}
\begin{enumerate}[resume]
\item По известным плотностям стадий получите формулу свёртки и выведите плотность \(f_T\) и функцию распределения \(F_T(t)=\mathbb P(T<t)\) \emph{для своего семейства}. Укажите область значений и граничные случаи.
\item Для порога \(\tau=1{,}55\,\mathbb E T\) вычислите \(q_0=\mathbb P(T>\tau)\). Рассмотрим два отдельных вмешательства (в каждом изменяется параметр только одной стадии):
\begin{itemize}
\item E: увеличить интенсивность выбранного перехода в \(4/3\) раза, так что \emph{среднее} время этой стадии уменьшится на 25\%;
\item U: уменьшить период выбранного приёмника до 75\% исходного.
\end{itemize}
Рассчитайте \(q_1,q_2\) для вмешательства в первый или второй модуль. Сравните \(q_0-q_1\) и \(q_0-q_2\) – \emph{абсолютные снижения риска}; при желании оцените и относительные снижения \(1-q_i/q_0\). До контрольного опыта зафиксируйте все три риска, ожидаемый эффект каждого изменения и объяснение различий через параметры и закон суммы. Проверьте оба прогноза на одной серии из 120 новых проходов: в компьютерной модели сравните исходный режим и оба вмешательства на одних и тех же случайных парах \(X,Y\). Объясните, почему наблюдаемые частоты могут отличаться от расчётных, а их порядок – от теоретического.
\end{enumerate}

\Needspace{14\baselineskip}\section{Большая партия: ЗБЧ и ЦПТ}
\begin{enumerate}[resume]
\item Найдите \(\mu=\mathbb E T\), \(\sigma^2=\operatorname{Var}T\). Для \(n=100\) независимых проходов приближённо рассчитайте с помощью ЦПТ \(\mathbb P(\overline T_n\le 1{,}04\mu)\). Сначала введите и подтвердите свои значения $\mu$ и $\sigma$, затем зафиксируйте численный прогноз и сравните его с 300 моделируемыми партиями. Режимы $n=5$ и $n=20$ являются дополнительными; обязательный результат получается для $n=100$. Объясните, как здесь проявляются ЗБЧ и ЦПТ. Среднее индивидуальное время партии \(\overline T_n\) \emph{не равно} времени завершения всей параллельной партии.
\end{enumerate}

\Needspace{4\baselineskip}\section{Что сдавать и как оценивается работа}
Отчёт должен быть \textbf{кратким, содержательным и проверяемым}. В нём должны быть:
\begin{itemize}[topsep=4pt,itemsep=2pt]
\item вывод по первым наблюдениям и объяснение роли датчика B;
\item законы двух стадий, предположения, совместная и условная плотности;
\item один вывод закона суммы через свёртку;
\item \(q_0,q_1,q_2\), абсолютные снижения риска и сравнение с контрольными частотами;
\item прогноз по ЦПТ и объяснение расхождения с результатами 300 партий.
\end{itemize}
Журнал JSON/CSV прикладывайте только по указанию преподавателя. Не пересказывайте нажатия кнопок, не прикладывайте весь код и не дублируйте экранные графики. Рисунки нужны лишь тогда, когда помогают обосновать вывод.

\begin{studentnote}
\textbf{Оценивание (всего 20 баллов).} Наблюдение и обоснование датчика B – 3; вероятностная модель и предположения – 4; совместное и условное распределения – 4; свёртка и сравнение снижения риска – 5; ЗБЧ и ЦПТ – 3; интерпретация результатов – 1.
\end{studentnote}
Разрешены Python, символьные вычисления и генеративный ИИ. Оценивается понимание модели, областей интегрирования, самостоятельная интерпретация и защита полученных результатов, а не объём написанного кода. Формальные статистические тесты не требуются.

\par\smallskip\noindent\textbf{На защите.} Объясните необходимость B, границы условной плотности и свёртки, общий порог и абсолютные снижения риска и применение ЦПТ к среднему времени.
\end{document}
